\maketitle

\section*{The programme's questions}

The programme reconstructs structure from relations, not spaces.
Point-spaces (the real line, Stone spaces, Hilbert spaces) are
outcomes of the reconstruction, not its starting material.  This
is a bet: that the relational description --- contexts, reference
frames, directed refinement --- is the right starting point, and
that the structures usually assumed (points, $\sigma$-algebras,
global states) are \emph{reconstructed} from it.

\medskip\noindent
Three central questions organize the reading:

\begin{enumerate}
  \item \textbf{What must be added to observation to get structure?}

    The persistent question across all phases:
    \begin{itemize}[nosep]
      \item Phase 1: observation $\to$ dynamics (semigroup)
      \item Phase 2--3: observation $\to$ geometry
        (curvature, metric, action)
      \item Phase 4--6: observation $\to$ probability (CE),
        realization (descent), value-definiteness (distributivity)
    \end{itemize}

    \smallskip\noindent
    The reading form: \emph{Where does someone identify the exact
    condition that forces a valuation to behave --- and what
    structure does that condition live in?}

  \item \textbf{Reconstruction without points.}

    When can observation be modeled relationally --- via contexts
    and reference frames --- rather than by presupposing an
    underlying set-theoretic structure of points?  Numerical
    identity (point-evaluation) is the degenerate case where all
    contexts contract to a single point.

    \smallskip\noindent
    Paper~I does this for the Boolean case: probability is
    reconstructed from directed refinement of finite Boolean
    algebras, without assuming a sample space.  The open question
    is whether this extends beyond the Boolean case, and if so
    what plays the role of ``directed refinement'' when the
    algebra is not distributive.

    \smallskip\noindent
    \emph{(Sharpened by Zafiris, Ch.~6 of Foundations of
    Relational Realism: the relational description is not a
    technical convenience but the starting point; point-spaces
    are what you reconstruct, not what you assume.)}

  \item \textbf{Local-to-global extension.}

    When does locally defined information (on each Boolean context)
    determine a global object?  What are the \emph{conditions on
    the site} --- not the individual valuations --- that control
    whether local data extend globally?

    \smallskip\noindent
    CE is one instance: local charges always glue finitely
    additively (Kolmogorov), but gluing to a $\sigma$-additive
    measure requires something extra --- and that ``something
    extra'' is not a sheaf condition on any known topology
    (Zafiris 2006, Biesel 2024: dictionary translation, not
    theorem).  Paper~II's EA/PR/VDR is another instance: the
    obstruction to global state extension in the non-Boolean case
    is non-distributivity rather than failure of $\sigma$-additivity.

    \smallskip\noindent
    The general form: both the Boolean and non-Boolean extension
    problems are local-to-global questions on different categories.
    What controls the answer is the structure of the site, not the
    properties of individual states.

    \smallskip\noindent
    \emph{(Sharpened by Zafiris, Ch.~6: the gluing conditions
    on overlapping Boolean reference frames are the structural
    content; the topology of the site is what does the work.)}
\end{enumerate}

% ==================================================================
\section{Caramello --- \emph{Theories, Sites, Toposes} (2018)}

\textbf{READ (2026-05-17).  Verdict: the bridge technique cannot
see CE.  $\sigma$-additivity is not expressible in finitary
geometric logic, so it cannot appear as a quotient in Caramello's
framework.}

\medskip\noindent
Read Chapters 1, 3--10.  The book develops the ``bridge technique'':
given a geometric theory $\mathbb{T}$ of presheaf type with
canonical Morita equivalence
\[
  \mathbf{Sh}(\text{f.p.}\mathbb{T}\text{-mod}(\mathbf{Set})^{op},
  J_{\mathrm{at}})
  \;\simeq\;
  \mathbf{Sh}(\mathcal{C}_{\mathbb{T}'}, J_{\mathbb{T}'}),
\]
transfer properties between the syntactic site
$(\mathcal{C}_{\mathbb{T}'}, J_{\mathbb{T}'})$ and the semantic
site $(\text{f.p.}\mathbb{T}\text{-mod}(\mathbf{Set})^{op},
J_{\mathrm{at}})$ to obtain new results.

\medskip\noindent
\textbf{The discriminating constraint.}  A ``quotient'' of a
geometric theory $\mathbb{T}$ means adding a \emph{geometric}
sequent --- built from finite conjunctions, small disjunctions,
existential quantification.  $\sigma$-additivity requires
\emph{countable} universal quantification and countable
conjunction: it is not a geometric axiom.  So $\sigma$-additivity
literally cannot be expressed as a quotient of any geometric theory
in Caramello's framework.  The bridge technique cannot see the
finite/countable boundary where CE lives --- for the same
structural reason that Deligne/Frot cannot: finitary
expressiveness.

\medskip\noindent
\textbf{The infinitary escape.}  At $\kappa = \omega_1$
(Esp\'{i}ndola's generalization), $\sigma$-additivity becomes
expressible in $\kappa$-geometric logic.  But we already
established (\S 8 above) that Property~T at $\omega_1$ on the CE
site reduces to classical conditions.  The bridge either can't see
CE (finitary) or reduces it to known content (infinitary).

\medskip\noindent
\textbf{What the book does provide:}
\begin{itemize}[nosep]
  \item Theorem~3.2.5 (Duality): bijection between quotients of
    $\mathbb{T}$ and subtoposes of $\mathbf{Set}[\mathbb{T}]$.
  \item Chapter~4: coHeyting algebra on subtoposes,
    Booleanization, DeMorganization.
  \item Chapter~6 (\S 6.1.1): canonical Morita equivalences for
    presheaf-type theories.
  \item Chapter~10: applications --- Fra\"{i}ss\'{e}'s theorem
    topos-theoretically, Galois representations.
  \item p.~337(b): the theory of decidable Boolean algebras is of
    presheaf type; f.p.\ models $=$ finite Boolean algebras;
    homogeneous models $=$ atomless Boolean algebras.  But this is
    the \emph{algebraic} theory --- no measures, charges, or
    $\sigma$-additivity.
\end{itemize}

\medskip\noindent
\textbf{No measure-theoretic applications.}  Caramello has no
published work applying the bridge technique to measure theory,
valuations, or probability.

\medskip\noindent
\textbf{Contact with programme:} None directly.  The bridge
technique operates within finitary geometric logic, which cannot
express $\sigma$-additivity.  Sixth instance of the pattern: CE is
invisible to every finitary framework tried (Zafiris, Biesel,
Frot/Deligne, Esp\'{i}ndola-reduced, Caramello).

\medskip\noindent
\textbf{Source:} Oxford University Press, 2018.
ISBN 978-0-19-875891-4.

% ==================================================================
\section{Vickers --- \emph{Topology via Logic} (1989)}

\textbf{READ (2026-05-17).  Verdict: provides the framework
(locales, sublocales, Stone duality) but not the application to
measures.  The question remains open --- requires later literature.}

\medskip\noindent
The book is a domain theory / logic textbook.  It develops:
\begin{itemize}[nosep]
  \item Locales (\S 5.4): points $=$ frame homomorphisms to
    $\mathbf{2}$; spatialization/localization functors.
  \item Sublocales (\S 6.2): $=$ nuclei on frames $=$ frame
    quotients.  The sublocale of ``realized points'' (principal
    ultrafilters) is well-defined in this framework.
  \item Spectral algebraic locales (\S\S 9.2--9.3): coherent
    presentations, compact coprime opens, Stone duality.
  \item Stone spaces (\S 9.5): spectral $+$ Hausdorff
    $\leftrightarrow$ Boolean algebra.  Clopens $=$ elements of the
    Boolean algebra.  Booleanization functor.  Patch topology.
  \item $\omega$-algebraicity (\S 9.4): second countability,
    $\omega$-chain completions.  The finite/countable boundary
    appears here --- finite locales are spectral algebraic
    (Prop.~9.4.1); second countable $\leftrightarrow$
    $\omega$-algebraic (Thm.~9.4.6).
\end{itemize}

\medskip\noindent
\textbf{What the book does NOT provide:}
\begin{itemize}[nosep]
  \item No theory of valuations/measures on locales.
  \item No discussion of $\sigma$-additivity or CE in
    locale-theoretic terms.
  \item No characterization of ``a valuation is supported on the
    sublocale of principal ultrafilters.''
\end{itemize}

\medskip\noindent
\textbf{Original question:} Is CE expressible as a frame-internal
property of a valuation?  \textbf{Answer: cannot be determined from
this book.}  The framework is here (locales, sublocales, Stone
duality for Boolean algebras), but the theory of valuations on
locales is post-1989: Heckmann (1996), Vickers (2008, ``A localic
theory of lower and upper integrals'').  The reading direction
remains open but now points to Vickers's later work on localic
measure theory, not this textbook.

\medskip\noindent
\textbf{Useful residue:}
\begin{itemize}[nosep]
  \item \S 9.2 Theorem~9.2.2: a frame is \emph{coherent} iff it
    has a coherent presentation iff $A \cong \mathrm{Idl}(K)$ for a
    distributive lattice $K$.  This is the frame-level notion
    behind Deligne/Frot.
  \item \S 9.4 Definition~9.4.3: second countability $=$ countable
    basis.  The step from finite to countable basis is exactly the
    step from coherent to $\omega$-algebraic --- parallel to the
    logical step from finitary to $L_{\omega_1,\omega}$
    (Esp\'{i}ndola).
  \item \S 6.2 Proposition~6.2.8: sublocales characterized by
    closure under all meets $+$ Heyting implication condition.
    The sublocale machinery is available for formulating the
    support condition --- what is missing is the valuation side.
\end{itemize}

\medskip\noindent
\textbf{Contact with programme.}  The framework is now in hand.
The open question is whether Vickers's later localic measure theory
(2011) provides the valuation-on-locale machinery needed to
formulate ``$\mu$ is supported on the sublocale of principal
ultrafilters'' as a frame-internal condition.  This would be the
next reading direction if the CE question is revisited.

\medskip\noindent
\textbf{Source:} Cambridge University Press, 1989.  ISBN 0\,521\,36062\,5.

% ==================================================================
\section{R\'{e}dei \& Summers --- ``Quantum Probability Theory''
(2006)}

\textbf{READ (2026-05-17).  Verdict: review paper, Paper~II
reference.}

\medskip\noindent
Reviews quantum probability in the von Neumann algebra framework.
Classical probability $=$ abelian von Neumann algebra; quantum $=$
nonabelian.  Normal states ($=$ $\sigma$-additive) characterized
by continuity.  The type I/II/III classification controls which
probability-theoretic phenomena appear.

\medskip\noindent
\textbf{Original question:} What is the obstruction to extending
states across non-Boolean joins?  \textbf{Answer:}
Noncommutativity itself.  The type classification refines this
but no single algebraic condition beyond ``non-distributive
$+$ $\dim \geq 3$'' is isolated.

\medskip\noindent
\textbf{Useful for Paper~II:} clean reference for the parallel
between classical and quantum probability in the operator algebra
framework.

\medskip\noindent
\textbf{Source:} arXiv:quant-ph/0601158.

% ==================================================================
\section{D\"{o}ring \& Isham --- ``What is a Thing?'' (2008)}

\textbf{READ (2026-05-17).  Verdict: Paper~II reference.
Daseinisation is descent-adjacent but not CE-relevant.}

\medskip\noindent
The topos is $\mathbf{Sets}^{\mathcal{V}(\mathcal{H})^{op}}$:
presheaves on the poset $\mathcal{V}(\mathcal{H})$ of commutative
subalgebras of $B(\mathcal{H})$.  The spectral presheaf $\underline{\Sigma}$
assigns to each Boolean context $V$ its Gelfand spectrum; it has
no global elements ($=$ Kochen--Specker).  Daseinisation (\S 5.2--5.3)
maps a projection $\hat{P}$ to a sub-object of $\underline{\Sigma}$
by approximating $\hat{P}$ within each Boolean context.

\medskip\noindent
\textbf{Original question:} Does daseinisation give descent from
Stone space to realization?  \textbf{Answer:} Structurally adjacent
but different.  Daseinisation goes \emph{up} (approximating
non-Boolean by Boolean); descent goes \emph{down} (Stone space to
realized points).  The failure of global sections IS Kochen--Specker,
reformulated topos-theoretically --- the non-Boolean analogue of CE
failure, but with a different obstruction mechanism
(noncommutativity vs.\ non-$\sigma$-additivity).

\medskip\noindent
\textbf{Contact with programme:} Paper~II directly.  Not
CE-relevant.

\medskip\noindent
\textbf{Source:} arXiv:0803.0417.

% ==================================================================
\section{Fremlin Vol.\ 3, Ch.\ 39 --- measurable algebras}

\textbf{READ (2026-05-17).  Verdict: answers the question
precisely.  Weak $(\sigma,\infty)$-distributivity is the algebraic
CE condition.}

\medskip\noindent
\textbf{391D Theorem} (Kantorovich--Vulikh--Pinsker 1950).  A
Boolean algebra $\mathfrak{A}$ is measurable iff:
\begin{enumerate}[nosep]
  \item $\mathfrak{A}$ is Dedekind $\sigma$-complete,
  \item $\mathfrak{A}$ is weakly $(\sigma,\infty)$-distributive,
  \item $\mathfrak{A}$ is chargeable (admits a strictly positive
    finitely additive functional).
\end{enumerate}

\medskip\noindent
\textbf{391J Theorem} (Kelley).  $\mathfrak{A}$ is chargeable iff
$\mathfrak{A} = \{0\}$ or $\mathfrak{A} \setminus \{0\}$ is a
countable union of sets with non-zero intersection numbers.

\medskip\noindent
\textbf{Original question:} Is there a known algebraic condition
forcing every charge to be $\sigma$-additive?  \textbf{Answer:
Yes --- weak $(\sigma,\infty)$-distributivity} (plus
$\sigma$-completeness).  KVP 1950 is the theorem.  This is NOT a
tautological restatement: weak $(\sigma,\infty)$-distributivity is
a genuine lattice condition on how countable infima interact with
arbitrary suprema.

\medskip\noindent
\textbf{Contact with programme.}  This is the algebraic answer to
Question~3 (local-to-global extension) for the Boolean case.  The
condition controlling whether local charges extend to a
$\sigma$-additive global measure is a property of the
\emph{algebra} (the site), not the charges (the sections) ---
exactly what the programme's framing predicts.  Weak
$(\sigma,\infty)$-distributivity is the algebraic face of CE.

\medskip\noindent
\textbf{Key citation for Paper~I:} KVP 1950 (via Fremlin 391D)
characterizes the exact algebraic boundary.  This is what
Howson's ``missing completeness theorem'' looks like in algebraic
language.

\medskip\noindent
\textbf{Source:} Fremlin, \emph{Measure Theory} Vol.~3, Ch.~39,
freely available at
\texttt{https://www1.essex.ac.uk/maths/people/fremlin/chap39.pdf}.

% ==================================================================
\section{Vickers --- ``A localic theory of lower and upper
integrals'' (2008)}

\textbf{Question while reading:} Can ``a valuation is supported on
the sublocale of principal ultrafilters'' be formulated as a
frame-internal condition on a localic valuation?  Does
$\sigma$-additivity of the valuation emerge as a locale-theoretic
property (e.g., preservation of directed joins, or a continuity
condition on the valuation map)?

\medskip\noindent
\textbf{Why.}  Vickers (1989) provides the framework (locales,
sublocales, Stone duality) but not the application to measures.
This later work develops valuations on locales constructively.
The programme's question --- is CE a frame-internal property? ---
requires this machinery.  If the support condition
``$\mu(\text{sublocale of realized points}) = 1$'' can be expressed
without reference to points, it would be the non-circular
characterization of coherence.

\medskip\noindent
\textbf{Contact with programme:} CE characterization directly.
This is the actual target that Vickers (1989) pointed toward.

\medskip\noindent
\textbf{Source:} \emph{Math.\ Logic Quarterly} 54(1), 109--123.

% ==================================================================
\section{Howson --- ``De Finetti, Countable Additivity, Consistency
and Coherence'' (2008)}

\textbf{READ (2026-05-17).  Verdict: identifies the same gap
(finite additivity vs.\ $\sigma$-additivity) but does not
formalize it.  TERMINOLOGICAL WARNING: Howson uses ``consistency''
and ``coherence'' as synonyms for de Finetti's criterion (finite
additivity / no Dutch book).  $\sigma$-additivity sits above both.
This is opposite to the programme's usage, where consistency $=$
finite additivity and coherence $=$ the stronger condition
including $\sigma$-additivity.}

\medskip\noindent
Howson clarifies de Finetti's position on countable additivity
and shows it is internally consistent.  Key findings:

\medskip\noindent
\textbf{The mistranslation.}  De Finetti's Italian \emph{coerenza}
$=$ ``consistency'' (\emph{consistenza}).  His 1972 book uniformly
uses ``consistency.''  The English translators (Kyburg) imposed
``coherence'' against de Finetti's own usage (p.~3).  The
programme's terminology aligns with de Finetti's original, not the
mistranslation.

\medskip\noindent
\textbf{Consistency $=$ finite additivity, precisely.}  De Finetti's
criterion: $P$ is consistent iff no finite combination of bets at
the $P$-rates guarantees a loss.  Characterizes exactly the
finitely additive probability functions.  Two properties parallel
first-order logic: \emph{absoluteness} (consistency depends only
on $P[E]$, not on which algebra $E$ sits in) and
\emph{compactness} (inconsistency detected by a finite subset).
Both fail for $\sigma$-additivity (pp.~8--9).

\medskip\noindent
\textbf{De Finetti's argument against the Dutch Book for
$\sigma$-additivity (\S\S 3--4).}  The infinite fair lottery
($p_n = 0$ for all $n$) is Dutch-Bookable only if postulate~(A)
--- ``a finite sum of fair bets is fair'' --- is extended to
countable sums.  De Finetti: this extension \emph{entails}
$\sigma$-additivity, making the Dutch Book argument circular.

\medskip\noindent
\textbf{Cox $\to$ finite additivity only (\S 6, p.~19).}  Cox's
three scale-free axioms $+$ propositional logic $\to$ finitely
additive probability.  ``Cox's result does not extend to endorsing
countable additivity.  There are infinitary propositional
languages\ldots\ $L_{\omega_1,\omega}$\ldots\ but there is no
natural extension of Cox's proof which would exploit that
additional structure.''

\medskip\noindent
\textbf{The open problem (p.~17).}  ``Something is nevertheless
missing, and that is a type of completeness theorem telling us that
the rules of probability extend to finite but not countable
additivity.''  This is precisely the gap that Frot and
Esp\'{i}ndola formalize: Deligne (finitary completeness) $=$
consistency; CE requires $\omega_1$-completeness.

\medskip\noindent
\textbf{Contact with programme.}  Howson (2008) is the primary
prior-art citation for the companion note.  He identifies the gap
(finite additivity has compactness, $\sigma$-additivity doesn't)
and states the open problem (p.~17: ``missing completeness
theorem'').  The companion note makes this precise via ultraproduct
proof.  The Frot/Esp\'{i}ndola boundary is the formal version of
Howson's philosophical observation: Deligne (finitary completeness)
$=$ consistency; CE requires something beyond
($\omega_1$-completeness).  Cox's axioms are the right starting
material for the programme: they produce finite additivity from
structural principles; $\sigma$-additivity requires something extra
that is not grounded the same way.

\medskip\noindent
\textbf{Five-traditions unification: DEAD (2026-05-18).}
Abramsky's contextuality obstruction is compact (finite covers);
the $\sigma$-additivity gap is non-compact (invisible to finite
tests).  Category error.  Removing Abramsky reduces the thesis to
Howson (2008) with fancier coordinates.  No new theorem.

\medskip\noindent
\textbf{Source:} \emph{BJPS} 59(1):1--23.

% ==================================================================
\section{Frot --- ``G\"{o}del's Completeness and Deligne's Theorem''
(2013)}

\textbf{READ (2026-05-17).  Verdict: confirms the dictionary,
places CE.}

\medskip\noindent
Proves Deligne's theorem (``coherent topoi have enough points'')
is equivalent to G\"{o}del's completeness for first-order logic.
The dictionary: models $=$ points of the topos (geometric morphisms
$\mathcal{E} \to \mathbf{Set}$); consistency $=$ non-triviality;
coherent theory $=$ finitary axioms $=$ coherent topos.

\medskip\noindent
\textbf{Original question:} Is CE's failure exactly ``the
classifying topos lacks enough points''?  \textbf{Answer: Not
exactly.}  CE is not expressible in coherent (finitary) logic
(confirming \L o\'{s}), so Deligne's theorem does not apply.  CE
lives \emph{above} the coherent level --- the guarantee of enough
points stops before reaching $\sigma$-additivity.  The right
framing: CE requires a completeness theorem for the
\emph{infinitary} fragment where $\sigma$-additivity lives.

\medskip\noindent
\textbf{What connects.}  Finite additivity IS coherent ---
Biesel's result (sheaf for finite covers) is the measure-theoretic
face of Deligne.  $\sigma$-additivity requires countable
operations, escaping coherence.  Frot makes this boundary precise:
it is the boundary between finitary and infinitary logic,
equivalently between coherent and non-coherent toposes.

\medskip\noindent
\textbf{Source:} arXiv:1309.0389.

% ==================================================================
\section{Esp\'{i}ndola --- ``Infinitary Generalizations of Deligne's
Completeness Theorem'' (2017; published \emph{JSL} 85(3), 2020)}

\textbf{READ (2026-05-17).  Verdict: places CE at $\kappa =
\omega_1$; identifies Property~$T$ as the structural condition.}

\medskip\noindent
Generalizes Deligne to $\kappa$-geometric logic (conjunctions and
disjunctions of ${<}\,\kappa$ formulas, existential quantification
over ${<}\,\kappa$ variables).  Key results:
\begin{itemize}[nosep]
  \item \textbf{Theorem~2.11.}  If $\kappa$ is regular with
    $\kappa^{<\kappa} = \kappa$, $\kappa$-geometric theories of
    cardinality ${\leq}\,\kappa$ are complete w.r.t.\
    $\mathbf{Set}$-valued models.
  \item \textbf{Corollary~3.1.}  Under same hypothesis, every
    $\kappa$-separable topos has enough $\kappa$-points.
  \item \textbf{Theorem~3.4.}  If $\kappa$ is weakly compact, every
    $\kappa$-coherent topos has enough $\kappa$-points.
  \item \textbf{Corollary~3.3.}  $\kappa$-coherent theories of
    cardinality ${\leq}\,\kappa$ are complete ($=$ Karp's
    completeness for $\mathcal{L}_{\kappa,\kappa}$).
\end{itemize}

\medskip\noindent
\textbf{Property~$T$} (Definition~2.2): an exactness condition on
the category --- every proper diagram (tree with bar) is completely
proper.  Combines strong distributivity $+$ dependent choice.  At
$\kappa = \omega$, property~$T$ is automatic (Remark~2.7),
recovering classical Deligne.  At $\kappa > \omega$ it is a genuine
restriction.  Rule~$T$ (p.~1150) is the syntactic counterpart.

\medskip\noindent
\textbf{Where CE lives.}  $\sigma$-additivity uses countable
universal quantification and countable conjunction $\to$
expressible in $L_{\omega_1,\omega}$ $\to$ relevant level is
$\kappa = \omega_1$.  At this level:
\begin{itemize}[nosep]
  \item Theorem~2.11 applies under CH
    ($\omega_1^{<\omega_1} = \omega_1$).
  \item Theorem~3.4 requires $\omega_1$ weakly compact (large
    cardinal axiom).
\end{itemize}

\medskip\noindent
\textbf{Reduction check (2026-05-17).}  Property~$T$ at
$\kappa = \omega_1$ on the site of finite Boolean subalgebras
reduces to $\sigma$-completeness of the underlying algebra $+$
uniqueness of the Carath\'{e}odory extension.  Both classical,
both already assumed in any CE setup.  No gap for a theorem.
Fifth instance of the pattern: CE-adjacent seeds die because
the topos/sheaf/logical machinery, applied to the specific site
where CE lives, restates classical conditions.

\medskip\noindent
\textbf{Status.}  The \emph{coordinates} are genuine (CE lives at
$\kappa = \omega_1$ in Esp\'{i}ndola's hierarchy, above Deligne
but within Karp).  Useful for exposition (Paper~I) and for the
Caramello reading direction (specifies what her machinery would
need to do).  Does not produce a seed.  The Caramello direction
remains open --- her Morita equivalence approach is structurally
different from ``apply Property~$T$ to this site.''

\medskip\noindent
\textbf{Source:} arXiv:1709.01967.

% ==================================================================
\section{Pt\'{a}k --- ``Exotic logics'' (1987)}

\textbf{READ (2026-05-17).  Verdict: populates Square~A bottom-right.
Obstruction is combinatorial (Greechie pasting), not from the center.}

\medskip\noindent
\textbf{The construction.}  An \emph{exotic logic} is a
$\sigma$-orthocomplete OMP that admits finitely additive states but no
$\sigma$-additive states.  Pt\'{a}k constructs these by
\emph{Greechie pasting}: start with a collection of Boolean algebras
(blocks), identify certain atoms across blocks so that the resulting
structure is an OMP, and show that any $\sigma$-additive state would
require the atoms of one block to sum to more than~1 (by the
identifications forcing double-counting).

\medskip\noindent
\textbf{Original question:} Is the obstruction purely from the center,
or is there a genuinely non-Boolean mechanism?  \textbf{Answer:}
genuinely non-Boolean.  The obstruction is combinatorial: the atom
identifications across blocks create contradictory partition-of-unity
constraints.  Each block individually admits $\sigma$-additive states
(it is Boolean), but the pasting forces any global state to satisfy
incompatible normalization conditions at the $\sigma$-additive level.
The center of the resulting OMP can be prescribed independently
(Theorem~3: any logic embeds into an exotic logic with given center).

\medskip\noindent
\textbf{Key result (Theorem~1).}  There exist exotic logics of
arbitrary cardinality.  The construction is explicit: the pasting
diagram and the atom identifications are given.

\medskip\noindent
\textbf{Contact with programme:} Square~A directly.  Confirms that the
OML extension problem (bottom-right cell) has a different character
from the Boolean extension problem (top-right, CE): the obstruction is
combinatorial/geometric (from the pasting), not analytic (from failure
of $\sigma$-additivity on a single Boolean algebra).  Gleason's theorem
is the exception --- it requires $\dim \geq 3$ and orthomodularity,
and even then only forces $\sigma$-additivity on specific algebras.

\medskip\noindent
\textbf{Source:} \emph{Colloquium Math.}\ 54(1), 1--7.

% ==================================================================
\section{Navara --- ``Regularity and $\sigma$-additivity of states''
(1992)}

\textbf{READ (2026-05-17).  Verdict: counterexample kills the
conjecture that regularity forces $\sigma$-additivity.  Not an
admissibility condition.}

\medskip\noindent
\textbf{The result.}  B\'{e}aver and Cook (1989) proved: on
\emph{projection lattices of von Neumann algebras}, every
$P$-regular state (= inner regular w.r.t.\ finite-dimensional
projections) is $\sigma$-additive.  Navara shows this does NOT
generalize: there exist $\sigma$-orthocomplete quantum logics
(OMPs) on which all states are $P$-regular yet some are not
$\sigma$-additive.

\medskip\noindent
\textbf{Original question:} Does regularity force $\sigma$-additivity
on general quantum logics?  Is this the admissibility condition for
the right column?  \textbf{Answer: No.}  The B\'{e}aver--Cook result
is special to von Neumann algebras (where the projection lattice has
additional structure: completeness, modularity in finite type,
type decomposition).  On general $\sigma$-orthocomplete OMPs,
regularity is strictly weaker than $\sigma$-additivity.

\medskip\noindent
\textbf{The construction.}  Navara modifies Pt\'{a}k's exotic-logic
construction: takes a Greechie pasting that admits only non-$\sigma$-additive
states, then verifies that all states on this logic are $P$-regular
(because the finite-dimensional sub-logics approximate the full logic
sufficiently well for inner regularity, but not for
$\sigma$-additivity).

\medskip\noindent
\textbf{Contact with programme:} Clarifies the bottom-right cell of
Square~A.  The OML extension problem does \emph{not} have a simple
admissibility condition analogous to regularity.  The gap between
finitely additive states and $\sigma$-additive states on OMPs is
not bridgeable by regularity alone --- the obstruction (Pt\'{a}k's
combinatorial pasting) persists even when all states are regular.
This confirms that the non-Boolean extension problem is
structurally harder than the Boolean one (where weak
$(\sigma,\infty)$-distributivity does the job via KVP~1950).

\medskip\noindent
\textbf{Source:} \emph{Proc.\ AMS} 115(2), 427--429.

% ==================================================================
\section{Shultz --- ``A characterization of state spaces of
orthomodular lattices'' (1974)}

\textbf{READ (2026-05-17).  Verdict: orthomodularity imposes NO
geometric constraint on state spaces; but $\sigma$-additive state
spaces CAN be geometrically degenerate.}

\medskip\noindent
\textbf{Main Theorem (p.~321).}  If $C$ is a compact convex subset
of a locally convex space, then $C \cong S(L)$ (affinely
homeomorphic to the state space) of some $\sigma$-orthomodular
lattice~$L$.  If $C$ is a rational polytope, $L$ can be chosen
finite.

\medskip\noindent
\textbf{Construction.}  Uses Greechie's $G$-spaces (the same
pasting technique as Pt\'{a}k's exotic logics): points $+$ frames
with loop constraints.  Lemma~1: any $G$-space yields a
$\sigma$-OML whose state space is affinely homeomorphic to the
weight space $W(X, \mathscr{F})$.  Lemmas~2--3: forcing techniques
(adjoining points, controlling weight values) to realize arbitrary
compact convex sets as weight spaces.

\medskip\noindent
\textbf{The negative result for $\sigma$-states (p.~327).}  The
converse is FALSE for $\sigma$-additive states.  Take $L$ $=$
Borel subsets of $\mathbb{R}$ modulo Lebesgue-null sets (a
$\sigma$-Boolean algebra).  Then $S_\sigma(L)$ has no extreme
points --- by Krein--Milman, $S_\sigma(L)$ is NOT affinely
equivalent to any compact convex subset of a locally convex space.
The $\sigma$-state space can be geometrically degenerate in ways
the full state space cannot.

\medskip\noindent
\textbf{Open question (p.~327).}  How to characterize $S(L)$ if
$L$ admits an order-determining set of states?  Shultz conjectures
there exist compact convex subsets not representable even with this
extra assumption.

\medskip\noindent
\textbf{Contact with programme:} Square~A, bottom row.
\begin{enumerate}[nosep]
  \item Orthomodularity alone imposes NO constraint on state-space
    geometry --- any compact convex set arises.  The obstruction to
    $\sigma$-additivity (Pt\'{a}k) cannot be read off from convex
    geometry.
  \item $S_\sigma(L)$ is structurally different from $S(L)$: it can
    lack extreme points entirely.  This is not just ``a smaller
    convex set'' --- it can fail to be a compact convex set at all
    in the relevant sense.
  \item Greechie pasting is the common tool behind both Shultz
    (state-space universality) and Pt\'{a}k (exotic logics with no
    $\sigma$-states).  The pasting is the fundamental technique for
    the non-Boolean column.
\end{enumerate}

\medskip\noindent
\textbf{Source:} \emph{J.\ Combin.\ Theory~A} 17(3), 317--328.
DOI: 10.1016/0097-3165(74)90096-X.

% ==================================================================
\section{Epperson \& Zafiris --- \emph{Foundations of Relational
Realism} (2013)}

\textbf{READ (2026-05-17).  Verdict: does not address CE.}

\medskip\noindent
Read Chapters 6, 8, 9, 10 of the book and Zafiris (2006)
\emph{J.\ Math.\ Phys.}\ 47, 092103 (the paper that puts measures
on the site).  The Grothendieck topology $J$ (epimorphic families)
does not grade covers by cardinality --- finite, countable, and
uncountable covers satisfy $J$ uniformly.  The finite/countable
boundary where CE lives is invisible.  States are finitely additive
throughout; $\sigma$-completeness is an axiom, not a sheaf
condition.  No cohomology is computed; the paper proves a
\emph{reconstruction} theorem (counit is iso), not an obstruction
theorem.  The $H^1$ unification with Abramsky--Brandenburger does
not appear.

\medskip\noindent
\textbf{Original question:} Does Zafiris's Boolean localization
give a sheaf-theoretic formulation of when local states glue to a
global $\sigma$-additive state?  \textbf{Answer: No.}

% ==================================================================
\section{Biesel --- ``Sheaves of Probability'' (2024)}

\textbf{READ (2026-05-17).  Verdict: confirms dictionary translation.}

\medskip\noindent
Proves measures form a sheaf for finite covers (Thm~8),
probability measures likewise (Thm~11).  Countable covers
explicitly excluded (Remark~7: uniform measures on $\{1,\ldots,n\}$
fail to glue to $\mathbb{N}$).  Section~5 states explicitly that
the results hold for merely finitely additive charges.  For
countable covers, $\sigma$-finiteness is needed --- standard
measure theory, no hidden structure.

\medskip\noindent
\textbf{Original question:} Is the step from finite to countable
covers exactly the CE boundary?  \textbf{Answer: Yes, but
tautologically.}  Finitely additive $=$ sheaf for finite covers;
$\sigma$-additive $=$ sheaf for countable covers.  This is a
restatement, not a characterization.

\medskip\noindent
\textbf{Useful residue:} Thms~8 and~11 are clean citations for
Paper~I's CE discussion.  Ref~[2] (Ross 2012, ``All roads lead to
violations of countable additivity,'' \emph{Phil.\ Studies})
relevant for Howson/de~Finetti direction.

\medskip\noindent
\textbf{Source:} arXiv:2401.01968.

% ==================================================================
\section*{Open leads: literature audit results (2026-05-17)}

\subsection*{Strategy D}

\textbf{Core question:} Does there exist a non-$\sigma$-complete,
non-atomic Boolean algebra admitting no $\sigma$-additive probability?

\textbf{Key structural clarification.}  On the \emph{clopen} algebra
of a Stone space, $\sigma$-additivity is vacuous by compactness:
there are no nontrivial countable disjoint families whose join is
also clopen.  The real Strategy~D question is about \emph{Radon
measures} on the Stone space (the Borel/Baire $\sigma$-algebra
generated by clopens).

\textbf{Literature findings:}
\begin{itemize}[nosep]
  \item \textbf{Plebanek (2024 survey):} most current secondary
    source on strictly positive measures on compact spaces.
    Covers Kelley condition, chain conditions, Fremlin's problem
    list.  Confirms the exact parameter regime (non-basically-disconnected,
    zero-dimensional, non-atomic) is open.
  \item \textbf{Argyros (1983):} constructs (under CH) a compact,
    totally disconnected, non-atomic, Radon-measure-free space ---
    but it is basically disconnected (corresponding Boolean algebra
    is $\sigma$-complete).  Wrong row of the hierarchy.
  \item No ZFC construction or obstruction is known for Strategy~D's
    exact conditions.  Likely independent.
\end{itemize}

\textbf{Verdict:} genuinely open, positioned precisely in the
landscape.  Next step would be consistency constructions
(diamond/CH) or reduction to known independence results.

\subsection*{OML extension problem}

\textbf{Core question:} What condition on a non-Boolean OML forces
every state to be $\sigma$-additive (or extendable to a
$\sigma$-additive measure on the McDonald--Bimb\'{o} dual)?

\textbf{The structural diagnosis.}  Pt\'{a}k--Pulmannov\'{a} (1994)
prove: if every unital subadditive measure on an OMP is a state,
the lattice must be Boolean.  This is the killer result ---
conditions strong enough to force $\sigma$-additivity destroy the
non-Boolean structure.

\textbf{Literature findings:}
\begin{itemize}[nosep]
  \item \textbf{Bunce--Wright (1992):} $\sigma$-additivity via
    operator-algebraic structure (JBW-algebras).  Positive but
    narrow: requires completeness close to von Neumann.
  \item \textbf{Chetcuti--Dvure\v{c}enskij (2003, 2005):} lattice
    effects algebras; positive results under strong completeness
    conditions.
  \item \textbf{Navara (1992):} regularity does NOT force
    $\sigma$-additivity on general OMPs.
  \item \textbf{Recent incremental work:} Vor\'{a}\v{c}ek--Pt\'{a}k
    (2023), Bure\v{s}ov\'{a}--Pt\'{a}k (2023), De~Simone--Navara
    (ongoing).  No breakthrough.
  \item \textbf{Wilce (2009), Handbook of Quantum Logic Ch.~24:}
    confirms Greechie pasting is the universal construction tool
    for exotic OML state spaces.
\end{itemize}

\textbf{Verdict:} genuinely open, structurally resistant.  The gap
between ``too weak'' and ``too strong'' appears robust.  Needs a new
idea (novel pasting technique, or categorical/topos bypass), not
more reading.

% ==================================================================
\section*{Context from literature audits (2026-05-15 to 2026-05-17)}

The consistency/coherence gap is real and recognized implicitly in
five traditions, but never unified under a single name:
\begin{itemize}[nosep]
  \item Philosophy of probability (de Finetti/Howson): consistency
    $=$ finite additivity; $\sigma$-additivity is beyond.  Howson
    (2008) confirms and defends de Finetti's position.
  \item Topos theory (Caramello/Johnstone): ``coherent theory'' $=$
    compactness-controlled; non-coherent $=$ beyond.  Frot (2013)
    proves Deligne $=$ G\"{o}del completeness.
  \item Infinitary logic (Esp\'{i}ndola 2020): CE lives at
    $\kappa = \omega_1$ in the $\kappa$-coherent hierarchy, above
    Deligne but within Karp.  Property~$T$ at $\omega_1$ reduces to
    classical conditions on the CE site --- no gap for a theorem.
  \item Quantum foundations (Abramsky et al.): contextuality as
    obstruction to global sections (sheaf cohomology).
  \item Imprecise probability (Walley): conglomerability $=$
    ``$\sigma$-coherence'' in all but name.
  \item Continuous model theory (Ben Yaacov): sidestep --- change the
    logic so $\sigma$-additivity becomes first-order.
\end{itemize}
No one has unified these.  The unification is the opportunity;
the risk is that it may be terminological packaging rather than
a theorem.

\medskip\noindent
The boundary is now precisely located from three angles:
\begin{itemize}[nosep]
  \item \textbf{Philosophical:} Howson --- Cox's axioms produce
    finite additivity only; no ``natural extension'' to
    $L_{\omega_1,\omega}$ gives $\sigma$-additivity.
  \item \textbf{Logical:} Frot --- Deligne (enough points for
    coherent topoi) $=$ finitary completeness $=$ consistency.
    $\sigma$-additivity escapes coherence.
  \item \textbf{Topos-theoretic:} Esp\'{i}ndola --- $\omega_1$-coherent
    completeness exists but Property~$T$ on the CE site collapses
    to classical conditions.  The machinery doesn't produce new
    content at this site.
\end{itemize}

\medskip\noindent
\textbf{Earlier findings (2026-05-15 to 2026-05-17):}
\begin{itemize}[nosep]
  \item The Simpson sublocale conjecture is \textbf{dead}: $\delta_U$
    counterexample kills $\Leftarrow$; compact Hausdorff makes
    $\Rightarrow$ trivially true.
  \item The sheaf/gluing approach (CE as descent) confirmed as
    \textbf{dictionary translation} by Zafiris (2006) and Biesel
    (2024): $\sigma$-additivity IS the sheaf condition for
    countable-partition topology by definition.
  \item Resolution-dependent charges / fiber defect: \textbf{dead}.
    Collapses to tail mass (Astbury 1981).  Interval-charge framework
    produces no theorem.
\end{itemize}

% ==================================================================
\section*{Openness as primitive (exploratory, 2026-05-19)}

The programme's central question --- what forces the transition from
finite to countable additivity --- discovers incompleteness as a
\emph{consequence}: the finitely additive framework is definite and
consistent; $\sigma$-additivity is the extra structure that fails to
be derivable from it.  A complementary question: are there frameworks
that take openness, indeterminacy, or incompleteness as \emph{primitive},
and if so, do they make contact with the FA/$\sigma$A boundary?

This section collects reading directions for that question.  Each
entry carries a specific discriminating question.  The question is
what determines whether the reference produces a Phase~1 seed or
gets parked.

\subsubsection*{Dzhafarov--Kujala, Contextuality by Default}

\noindent\textbf{Primary:} E.~N.~Dzhafarov and J.~V.~Kujala,
``Contextuality-by-Default 2.0,'' \emph{Lecture Notes in Computer
Science} 10106 (2017), 16--32.  See also Dzhafarov and Kujala,
``Context-content systems of random variables,'' \emph{Philosophical
Transactions of the Royal Society A} 375 (2017), 20160389.

\medskip\noindent
CbD starts from the position that random variables measured in
different contexts \emph{lack a joint distribution primitively}.
A joint distribution (coupling) is an achievement, not a given.
Contextuality is then the failure of a maximally connected coupling
to exist.  The framework is probabilistic, measure-theoretic, and
has a precise notion of ``obstruction to global coherence.''

\medskip\noindent\textbf{Discriminating question:} Does CbD's
coupling-existence criterion have an analogue for the FA $\to$
$\sigma$A extension problem?  Concretely: can $\sigma$-additivity
be characterized as the existence of a coupling across refinement
contexts --- i.e., the existence of a joint distribution on
$\bigcup_k \mathcal{G}_k$ compatible with each $\mu|_{\mathcal{G}_k}$?
If yes, this would give a non-trivial structural characterization.
If the analogy is only verbal (``obstruction to coherence'' in
both cases but different mathematical content), park it.

\subsubsection*{Ellerman, Partition Logic and Information --- PARKED}

\noindent\textbf{Primary:} D.~Ellerman, ``An Introduction to
Partition Logic,'' \emph{Logic Journal of the IGPL} 22 (2014),
94--125.  See also Ellerman, ``The Logic of Partitions,''
\emph{Review of Symbolic Logic} 3 (2010), 287--350; and
\emph{New Foundations for Information Theory}, Springer (2021).

\medskip\noindent
\textbf{Status: PARKED (2026-05-19).}  Read in full.
Ellerman builds a complete dual logic where partitions (equivalence
relations) play the role that subsets play in Boolean logic.
Distinctions (``dits'') dualize elements; the partition lattice
$\Pi(U)$ is non-distributive; each partition $\pi$ generates a
Boolean core $\mathcal{B}_\pi \cong \mathcal{P}(\pi_{ns})$.
``Logical entropy'' $h(\pi) = \sum p_i(1-p_i)$ is the normalized
dit-counting measure.

\medskip\noindent
\textbf{Verdict on discriminating question:}  The dit/bit translation
\emph{is} relabeling entropy.  Logical entropy $h(\pi) = 1 -
\sum p_i^2 = 1 - \operatorname{Coll}_\mu(\mathcal{G})$, i.e., one
minus the collision probability already used in the programme's
observational-resolution notes.  The partition refinement order
recovers the cylinder-algebra filtration.  The paper stays entirely
in the finite setting --- no measures, no charges, no limit behavior,
no extension problems.  No theorem in partition logic translates to
new content about the FA/$\sigma$A boundary.

\medskip\noindent
\textbf{Possible re-scoping (not pursued):} The non-distributivity
of $\Pi(U)$ and the Boolean core construction could speak to the
OML extension problem (where distributivity fails), but this would
be a different question from the one that motivated the entry.

\subsubsection*{Walley, Imprecise Probability}

\noindent\textbf{Primary:} P.~Walley, \emph{Statistical Reasoning
with Imprecise Probabilities}, Chapman \& Hall (1991).  For a concise
treatment, see M.~C.~M.~Troffaes and G.~de Cooman, \emph{Lower
Previsions}, Wiley (2014).

\medskip\noindent
Walley takes as primitive a \emph{set} of probability measures (a
credal set), not a single measure.  Sharp (precise) probability is
the special case where the credal set is a singleton.  The framework
has a natural hierarchy: upper/lower previsions $\to$ coherent
previsions $\to$ precise previsions $\to$ $\sigma$-additive previsions.
Each inclusion is strict.  The theory of ``natural extension'' ---
the smallest coherent prevision dominating a given assessment --- is
formally analogous to the problem of extending a finitely additive
charge.

\medskip\noindent\textbf{Discriminating question:} Is there a
Walley-type framing where the set of $\sigma$-additive extensions of
a finitely additive measure on a subalgebra is itself the primitive
object, and $\sigma$-additivity $=$ singleton extension set?  If so,
does this give a useful topology or duality on the extension sets
(connecting to Choquet theory, Phelps simplex structure)?  If the
Walley framework only redescribes the same extension problem in
different language without new theorems, park it.

\subsubsection*{Abramsky--Brandenburger, Sheaf-Theoretic Formalism}

\noindent\textbf{Primary:} S.~Abramsky and A.~Brandenburger, ``The
Sheaf-Theoretic Structure of Non-Locality and Contextuality,''
\emph{New Journal of Physics} 13 (2011), 113036.

\medskip\noindent
\textbf{Re-scoped.}  The contextuality application was previously
killed: the contextuality obstruction is compact (witnessed by finite
covers), while the $\sigma$-additivity gap is non-compact (invisible
to any finite test).  These are different phenomena.

What survives is the \emph{formalism}: a presheaf of local sections
(compatible local data), a sheaf condition (existence of a global
section), and a cohomological obstruction ($H^1$ or higher) measuring
the failure.  The question is whether this formal apparatus can be
applied to the Yosida-Hewitt decomposition, where the ``local
sections'' are the restrictions of a finitely additive charge to
finite subalgebras and the ``global section'' is a $\sigma$-additive
extension.

\medskip\noindent\textbf{Discriminating question:} Can the purely
finitely additive part $\mu_p$ of the Yosida-Hewitt decomposition be
expressed as a cohomological obstruction in the sense of
Abramsky-Brandenburger --- i.e., as a non-trivial class in $H^1$ of
a presheaf on a suitable category of finite subalgebras?  If so, does
the cohomological degree correspond to anything meaningful (e.g.,
``how far'' from $\sigma$-additivity)?  The key test is whether the
compact/non-compact mismatch is fatal to the formalism, not just the
application: the Yosida-Hewitt obstruction lives at countable limits,
and \v{C}ech cohomology of a Stone space is trivial for compact
coefficients.  If the cohomology is always zero, the formalism
doesn't apply.

\subsubsection*{Sambin, Positive Topology}

\noindent\textbf{Primary:} G.~Sambin, ``Some points in formal
topology,'' \emph{Theoretical Computer Science} 305 (2003), 347--408.
See also Sambin, \emph{Positive Topology: Completing the Foundational
Trilogy} (book in preparation, drafts circulating since 2020).

\medskip\noindent
Sambin's positive topology has two primitives: a \emph{cover relation}
(when does a collection of opens cover an open?) and a
\emph{positivity predicate} (which opens are ``inhabited'' ---
non-trivially instantiated).  Points are \emph{derived}, not
primitive: a point is a filter satisfying both cover and positivity
constraints.  The framework is constructive and works without
classical logic or the axiom of choice.

\medskip\noindent\textbf{Discriminating question:} Can the
FA/$\sigma$A transition be viewed as a completion in Sambin's sense,
where finitely additive charges live on a ``basic'' cover and
$\sigma$-additivity is the positivity predicate forcing concentration
on derived points?  Concretely: for a Boolean algebra $B$, is the
$\sigma$-additive part $\mu_c$ of a charge characterized by a
positivity condition on the Stone space (``every open containing a
point of $\operatorname{supp}(\mu_c)$ is positive''), and if so,
does this give a constructive proof of the Yosida-Hewitt
decomposition?  If the framework only redescribes Stone duality in
constructive language, park it.

\subsubsection*{Supporting philosophical references}

The following are not expected to produce Phase~1 seeds directly but
may sharpen the question of what ``starting from openness'' means.

\begin{itemize}[nosep]
  \item \textbf{Gisin (2021):} N.~Gisin, ``Indeterminism in
    Physics, Classical Chaos and Bohmian Mechanics,'' \emph{Erkenntnis}
    86 (2021), 1553--1590.  Argues that classical determinism is an
    artifact of completed infinities in the real-number continuum;
    intuitionistic reals leave room for genuine indeterminacy.
    \emph{Relevance:} if $\sigma$-additivity depends on completed
    countable operations, does an intuitionistic treatment change the
    extension boundary?
  \item \textbf{Linnebo--Shapiro (2019):} {\O}.~Linnebo and S.~Shapiro,
    ``Actual and Potential Infinity,'' \emph{No\^{u}s} 53 (2019),
    160--191.  Distinguishes potential from actual infinity with formal
    precision.  See also Hamkins's set-theoretic potentialism.
    \emph{Relevance:} the FA/$\sigma$A boundary is exactly where
    potential infinity (finite additivity: every finite test passes)
    meets actual infinity ($\sigma$-additivity: the countable limit
    exists).  Does the potentialist framework produce a formal
    distinction beyond the known one?
\end{itemize}

% ==================================================================
\section*{How to use this document}

When reading, annotate in the margins or on a separate sheet:
\begin{enumerate}[nosep]
  \item What theorem or definition surprised you.
  \item What connected to the programme.
  \item What contradicted the knowledge map.
  \item What suggested a concrete Phase~1 seed note.
\end{enumerate}

\noindent
When a reading direction produces a concrete claim worth auditing,
extract it to \texttt{notes/unsorted/} and run Phase~2.

